\chapter{Resum}
Aquest treball de recerca aborda un problema real de mobilitat a l’àrea metropolitana de Barcelona:
la manca de connexió entre els serveis de bicicletes compartides Bicing (Barcelona) i AMBici (municipis veïns).
L’objectiu principal ha estat desenvolupar un programa en Python capaç d’integrar les dades obertes d’ambdós serveis
mitjançant les seves APIs públiques, identificar punts de transbord i generar indicadors útils com el “canvi ràpid”.

El projecte també incorpora visualitzacions gràfiques, exportació de resultats a CSV i una aplicació web interactiva creada amb Streamlit,
que facilita l’accés als resultats sense coneixements tècnics previs.

Aquest treball combina aspectes tècnics (Python, Git, Linux, LaTeX) amb una dimensió social,
ja que fomenta la mobilitat sostenible i aporta una solució innovadora a un repte quotidià dels usuaris.

\chapter{Abstract}
This research project addresses a real mobility issue in the metropolitan area of Barcelona:
the lack of integration between the bike-sharing services Bicing (Barcelona) and AMBici (neighboring municipalities).
The main objective was to develop a Python program capable of integrating open data from both services through their public APIs,
identifying transfer points, and generating useful indicators such as the “quick change”.

The project also includes graphical visualizations, CSV export of results, and an interactive web application built with Streamlit,
which makes it easier to access the results without requiring technical knowledge.

This work combines technical aspects (Python, Git, Linux, LaTeX) with a social dimension,
since it promotes sustainable mobility and provides an innovative solution to a daily challenge faced by users.
